% =============================================================================
% Lecture 14: Supervisory Intervention
% Physical AI: Machine Learning Systems That Sense and Act (Volume IV), Chapter 14
% =============================================================================
% Thesis: a takeover is a control action with its own deadline, a single writer
% per actuator, a smooth transfer, and a record.
% Running case: the warehouse mobile manipulator (the running machine) at the
% rack end, 1.10 m ahead, at the 1.3 m/s aisle speed, handed to a remote assistant.
% Hook and synthesis: California DMV disengagement filings (SI 14.1).
%
% Build: make deck V=4 D=14_intervention   (figures: images/make_figs.py,
%        numbers: images/calcs.py)
% Plan:  lecture_map.yml in this folder.
% =============================================================================

\documentclass[aspectratio=169, 12pt]{beamer}
\usepackage{../../assets/beamerthememlsys}
\mlsyssetup{
  volume={Volume IV},
  chapter={Chapter 14},
  chaptertitle={Supervisory Intervention},
  layer=governance,
  logo={../../assets/img/logo-mlsysbook.png},
  instlogo={../../assets/img/logo-harvard.png},
}
\usepackage[T1]{fontenc}
\usepackage[scaled=0.9]{helvet}
\usepackage{courier}
\renewcommand{\familydefault}{\sfdefault}
\usepackage{booktabs,amsmath}
\graphicspath{{images/}}
\setcounter{mlsystotalsections}{4}

\title{Supervisory Intervention}
\author{Vijay Janapa Reddi}
\institute{Harvard University}
\date{}

% Diagram role colors (Volume IV palette), for text beside figures
\definecolor{rolegreen}{HTML}{047857}
\definecolor{rolered}{HTML}{991B1B}
\definecolor{roleblue}{HTML}{1D4ED8}
\definecolor{roleamber}{HTML}{B45309}
% A large answer choice for predict and concept-question slides
\newcommand{\choice}[2]{\textbf{\textcolor{mlsysaccent}{#1}}\enspace #2}
\newcommand{\answerat}[3]{\alt<#1->{\colorbox{mlsysaccent!15}{\choice{#2}{#3}}}{\choice{#2}{#3}}}
% A goal or takeaway line with its text-free thumbnail
\newcommand{\thumbline}[2]{%
  \begin{minipage}[c]{2.6cm}\centering\includegraphics[width=2.5cm,height=0.95cm,keepaspectratio]{#1}\end{minipage}%
  \hspace{0.3cm}\begin{minipage}[c]{0.74\textwidth}\small #2\end{minipage}}
% One line of a worked example, shown from overlay #1 on
\newcommand{\wstep}[2]{\uncover<#1->{#2}\\[0.2cm]}
% An exact excerpt of a public file, log, or code
\newenvironment{logblock}[1][7]{\par\ttfamily\fontsize{#1}{8.6}\selectfont\raggedright\setlength{\parskip}{0pt}}{\par}
\newcommand{\logrow}[2]{\makebox[\linewidth][l]{#1\hfill #2}\\}

\begin{document}

% =============================================================================
% TITLE AND HOOK
% =============================================================================
\mlsysnote{[0.5 min] [CORE] Read the thesis aloud once. This lecture opens Part IV, Governing the Machine, whose law is the fourth law, permission must be earned. Lecture 1 named all four laws and said that only the first is a physical law; the fourth is a design principle the book chooses to elevate. Today it applies for the first time to a person. Core minutes: about 57.5 of the 71 in these notes. FLEX: notes tagged [OPTIONAL] can be cut; cut first the DMV file excerpt, the session waterfall, the pallet, and the flight-display photo.}
\mlsystitle{Supervisory Intervention}{A takeover is a control action with its own deadline, a single writer per actuator, a smooth transfer, and a record.}{cover_intervention_blueprint.png}

\begin{frame}{Predict: will takeovers die out by 2025?}
  \mlsysactivity{Predict}{2 min}
  \centering
  \mlsyssteps[width=0.82\textwidth,height=4.6cm,keepaspectratio]{dmv_step}{2}\\[0.05cm]
  {\small One fleet's California filings: miles driven per takeover, \alt<2->{2015 to 2025; green is driverless testing.}{2015 to 2020.}}
  \mlsyscredit{California DMV autonomous vehicle disengagement reports, Waymo annual filings (via the MLSysIM registry); Supervisory Intervention, \S14.1, figure 1.}
  \note{[2.5 min] LINK: Part III built a permission path whose every delay belongs to the machine. Today a person enters the loop. NARRATE: California requires every company testing autonomous vehicles on its public roads to report its miles and its disengagements, the times control passed to a person. Step 1 shows one long-running fleet from 2015 to 2020: from 1,244 to 29,945 miles per takeover. ENGAGE: ask the room to predict 2021 to 2025. Will the curve keep rising until people are no longer needed? Take a show of hands. Step 2: the rate fell 3.8 times from 2020 to 2021 while the miles grew 3.7 times, and the latest year still has 174 takeovers at 19,234 miles each. Driverless testing (green) still reports one every 71,815 miles. WARN: students read the curve as a capability score. The rate depends on where the fleet drove and on what the company chose to count; the DMV says the reports are not designed for comparison across companies. FLEX: [CORE]. The synthesis returns to this plot.}
\end{frame}

\begin{frame}{Each filed takeover records who acted and why}
  \begin{logblock}[7.2]
    DATE~~~~~~~INITIATED BY~~LOCATION~~~DESCRIPTION OF FACTS CAUSING DISENGAGEMENT\\[0.05cm]
    11/28/2024~Test Driver~~~Street~~~~~Disengage for unwanted maneuver of the vehicle\\
    ~~~~~~~~~~~~~~~~~~~~~~~~~~~~~~~~~~~~that was undesirable under the circumstances\\[0.05cm]
    10/8/2024~~Test Driver~~~Freeway~~~~Disengage for a perception discrepancy for which\\
    ~~~~~~~~~~~~~~~~~~~~~~~~~~~~~~~~~~~~a component of the vehicle's perception system\\
    ~~~~~~~~~~~~~~~~~~~~~~~~~~~~~~~~~~~~failed to detect an object correctly\\[0.05cm]
    4/17/2024~~AV System~~~~~Freeway~~~~Disengage for unwanted maneuver of the vehicle\\
    ~~~~~~~~~~~~~~~~~~~~~~~~~~~~~~~~~~~~that was undesirable under the circumstances\\
  \end{logblock}
  \vspace{0.3cm}
  {\small Waymo, 2024 report: \textbf{244} rows, \textbf{230} begun by the test driver and \textbf{14} by the automated system. No column records how long the takeover took or who held each control afterward.}
  \mlsyscredit{California DMV, 2024 Autonomous Vehicle Disengagement Reports (CSV), three Waymo LLC rows, columns abridged; archived at web.archive.org, 13 Feb 2025.}
  \note{[1.5 min] NARRATE: these are real rows from the DMV's 2024 file, with the permit, VIN, and driverless-capable columns dropped to fit. Each row names a date, who began the disengagement, the road type, and a reason chosen from a short list. Of Waymo's 244 rows, 230 were begun by the test driver and 14 by the automated system (94.3 percent by the driver, derived). ENGAGE: if a robot struck a rack during one of these takeovers, which question could this file answer? (Expected: who began it and roughly why. It cannot say how long the person took, which controls moved, or what the machine did meanwhile.) WARN: a takeover count measures how often the exchange happened, not how well either party performed it. The rest of the lecture builds the record that can. FLEX: [OPTIONAL]. Cut and say one sentence about what the file holds.}
\end{frame}

\begin{frame}{Recall: how much spare does the aisle run keep?}
  \mlsysactivity{Retrieval}{1.5 min}
  \begin{columns}[c]
    \begin{column}{0.55\textwidth}
      {\Large\bfseries Close your notes.}\\[0.2cm]
      {\normalsize At 1.3~m/s toward the rack end 1.10~m ahead, how much clearance does the finished budget leave?}\\[0.3cm]
      \uncover<2->{\small \textbf{131.0~mm}, which is \textbf{100.8~ms} of extra delay at 1.3~m/s.\\[0.1cm]
      After the 63.3~ms dispatch-to-admission allocation: 48.7~mm, or \textbf{\textcolor{mlsysaccent}{37.5~ms}}.}
    \end{column}
    \begin{column}{0.41\textwidth}
      \centering\safeimg[width=\textwidth,height=2.2cm,keepaspectratio]{thumb_tax_step1.pdf}\\[0.05cm]
      {\footnotesize Allowances, machine delay, then the jerk-limited stop.}\\[0.2cm]
      \uncover<3->{\begin{mlsyscard}{mlsysaccent}\small\textbf{Today:} a person needs about 605~ms to take over.\end{mlsyscard}}
    \end{column}
  \end{columns}
  \mlsyscredit{Lecture 12, Safety Enforcement (finished budget, Napkin Math 12.1); Lecture 13, Workload Placement (dispatch-to-admission allocation). Values derived.}
  \note{[1.5 min] NARRATE: give 45 seconds. Step 2: Lecture 12 finished the stopping budget at the 1.3 m/s aisle speed. Its rows are 173.68 mm of travel during the 133.6 ms machine delay, a 605.31 mm jerk-limited stop, and 190 mm of fixed allowances, 968.99 mm in all, leaving 131.0 mm. At 1.3 m/s that spare is 100.8 ms (derived). Lecture 13 allocated 63.3 ms between dispatch and admission; charging it leaves 48.7 mm, or 37.5 ms. Step 3: ask the room to hold 100.8 ms in mind; today a human response of 605 ms arrives. ENGAGE: about how many times the spare is a 605 ms wait? (Expected: about six times.) FLEX: [CORE].}
\end{frame}

\begin{frame}{Why can a takeover cause the collision it should prevent?}
  \thumbline{thumb_strip_step3.pdf}{Price a takeover as travel before braking, and solve the budget for the speed and the deadline it allows.}\\[0.2cm]
  \thumbline{thumb_fsm_step3.pdf}{Bound every proposer with a lease, and move ownership so one source holds each channel per tick.}\\[0.2cm]
  \thumbline{thumb_quintic_step2.pdf}{Size a smooth command transfer and the clearance it costs.}\\[0.2cm]
  \thumbline{thumb_masks_step1.pdf}{Name the mode confusion a split handover invites, and log who held each channel so the record can be trusted and learned from.}\\[0.25cm]
  {\footnotesize\textcolor{midgray}{Fourth law: permission must be earned.\quad New record: the authority record.}}
  \note{[1.5 min] NARRATE: the driving question is the title, and it is the chapter's own. Four goals, one per segment; each thumbnail is the shape of the figure that teaches it. The running machine is the warehouse mobile manipulator, a wheeled base with an arm, driving toward a rack end where a person may step out. When a pallet blocks the aisle, it asks a remote assistant to take over. Most of its values are labeled illustrative or derived in the book; each segment pairs them with published measurements: driver take-over studies, safety standards, and an accident report. Pathways: machine learning students meet human-factors numbers and state machines; embedded students will recognize leases and heartbeats; control students will recognize bumpless transfer from process control. FLEX: [CORE].}
\end{frame}

\begin{frame}{A takeover moves the right to propose; the check stays}
  \centering
  \safeimg[width=0.9\textwidth,height=5.1cm,keepaspectratio]{propose_admit.pdf}\\[0.05cm]
  {\small Token: the right to propose on a channel. MCU: the safety microcontroller.}
  \mlsyscredit{Supervisory Intervention, \S14.2, table 1 and Definition 14.1.}
  \note{[2.5 min] LINK: Lecture 1 drew the proposal boundary: above it a component may only propose; below it the permission path admits. Lecture 12 built the enforcer on the safety microcontroller (Lecture 4). NARRATE: left, enforcement. The learned policy keeps proposing, and the enforcer changes a proposal that would break a limit. Right, intervention. The policy's proposal token is revoked on named channels and an authenticated operator receives it, but every operator command passes the same enforcer. Say the chapter's line: enforcement changes a proposal and leaves the proposer in place; intervention changes the proposer and leaves the check in place. ENGAGE: a startled operator commands a turn that would roll the chassis. What happens? (Expected: the enforcer clamps or rejects it, exactly as for the policy.) WARN: students assume a takeover hands the operator the actuators. The operator gains only the right to propose, and the machine still needs a local stop that does not wait for the operator. The third law, proposal is not permission, binds people and models alike. FLEX: [CORE].}
\end{frame}

% =============================================================================
% SECTION: DEADLINE
% =============================================================================
\mlsysnote{[0.5 min] [CORE] First the price of waiting for a person, in millimeters and in the speed the machine may keep while it asks.}
\mlsysdivider{How long can a moving machine wait for a person?}{Supervisory Intervention, \S14.1 and \S14.2}
\section{Deadline}

\begin{frame}{Remote operators see the machine only through video}
  \begin{columns}[c]
    \begin{column}{0.55\textwidth}
      \centering\safeimg[width=\textwidth,height=4.4cm,keepaspectratio]{rov-control-station.jpg}\\[0.03cm]
      {\footnotesize A pilot station for a remotely operated vehicle (ROV) in a ship's lab}
    \end{column}
    \begin{column}{0.42\textwidth}
      \centering\safeimg[width=\textwidth,height=3.0cm,keepaspectratio]{iss-robotics-workstation.jpg}\\[0.03cm]
      {\footnotesize The robotics workstation on the International Space Station, which drives the Canadarm2 arm}\\[0.2cm]
      {\small Every action crosses a camera, a link, a person, and a hand controller.}
    \end{column}
  \end{columns}
  \mlsyscredit{Photos: Dee Emrich, WHOI / Ocean Observatories Initiative, Wikimedia Commons; NASA, public domain, Wikimedia Commons.}
  \note{[1 min] NARRATE: two real remote-operation stations. Left, an ROV pilot in a darkened lab with video and telemetry screens and hand controllers. Right, an astronaut at the station's robotics workstation. A warehouse remote assistant works the same way: the scene arrives as video, the person decides, and a command goes back over a network. ENGAGE: list the delays between the event in the aisle and the first changed command. (Expected: camera and video encoding, transport, the person seeing, deciding, moving the hand, and the command's trip back.) WARN: faster networks trim the transport terms but cannot shorten the person's terms. FLEX: [OPTIONAL]. Cut and move to the timing strip.}
\end{frame}

\begin{frame}{An in-loop takeover spends 605 ms before any command}
  \centering
  \mlsyssteps[width=0.95\textwidth,height=5.0cm,keepaspectratio]{strip_step}{3}\\[0.05cm]
  {\small In-loop: hands on the controls, expecting the event. $T_\Sigma = t_{\text{perceive}} + t_{\text{decide}} + t_{\text{act}} + t_{\text{console}}$}
  \mlsyscredit{Supervisory Intervention, \S14.2.2 (eq.~1, figure 2); Card, Moran, and Newell 1983 (slow-bound cycles); Schimpe et al. 2022 (104 ms, wired link).}
  \note{[2.5 min] LINK: the spare from the recall slide is 131 mm, or 100.8 ms. NARRATE: Step 1: the scene reaches the operator after a video glass-to-glass latency of 104 ms, measured on a wired teleoperation link (Schimpe et al. 2022). The spare is gone at 100.8 ms, before the operator has seen anything. Step 2: the Model Human Processor prices a simple reaction as one perceptual, one cognitive, and one motor cycle; the design takes the slow end of each range: 200, 170, and 100 ms. Step 3: the four-phase handshake (this lecture's second segment) adds 31 ms, so the takeover commits at 605 ms. At 1.3 m/s that is 786.5 mm of travel, 6.0 times the spare (derived). ENGAGE: which term would a faster network shrink? (Expected: only video transport and the handshake, 135 ms of the 605.) WARN: the 104 ms was measured on a wired link; a warehouse wireless link must be measured on site. FLEX: [CORE].}
\end{frame}

\begin{frame}{Predict: does the stop still fit after a takeover?}
  \mlsysactivity{Predict}{2 min}
  \centering
  \mlsyssteps[width=0.92\textwidth,height=4.6cm,keepaspectratio]{tax_step}{2}\\[0.05cm]
  {\small At 1.3~m/s. Orange is the wait for the person; dark is the jerk-limited stop.}
  \mlsyscredit{Supervisory Intervention, Napkin Math 14.1 (chosen, illustrative, and derived values of the running machine; human terms from cited studies).}
  \note{[3 min] ENGAGE: show step 1, the machine alone, 968.99 mm of the 1,100 mm. Ask: insert the 605 ms in-loop takeover before braking. Fits or overruns, and by how much? Vote. NARRATE: Step 2: the in-loop takeover adds 786.5 mm and the stop reaches 1,755.49 mm, 655.5 mm past the rack end. The out-of-loop takeover, an assistant who was watching several machines, takes 4.5 s by the design value of the next slide; it adds 5,804.5 mm and the stop reaches 6,773.49 mm, 5.67 m past the line. All derived. Point out that the dark braking segment is identical in every bar. ENGAGE: would a stronger brake fix row two? (Expected: no; the waiting alone, 786.5 mm, exceeds the 131 mm spare six times.) WARN: students reach for better brakes. The waiting term, not the braking term, decides the verdict. FLEX: [CORE].}
\end{frame}

\begin{frame}{Take-over times run long: the 90th percentile is 4.33 s}
  \centering
  \safeimg[width=0.88\textwidth,height=4.7cm,keepaspectratio]{zhang_hist.pdf}\\[0.05cm]
  {\small 520 conditions in 129 driving studies. Under each line, the request speed that delay allows.}
  \mlsyscredit{Zhang, de Winter, Varotto, Happee, and Martens, Transportation Research Part F 64 (2019), supplementary data (CC0), replotted; ceilings from Supervisory Intervention, table 2 (derived).}
  \note{[2 min] NARRATE: this is real data, the 520 condition means of a meta-analysis of 129 driver take-over studies, replotted from the authors' published file. A driver who was monitoring automation and gets a take-over request needs on average 2.72 s; the median condition is 2.37 s; the 90th percentile is 4.33 s; the 99th is 7.0 s, and one condition took 19.8 s. Under each line is the speed ceiling the running machine's budget allows with that delay plus the 135 ms console term: 0.32, 0.29, 0.19, and 0.12 m/s (derived, table 2). The design uses the 90th percentile, so the out-of-loop takeover is 4.33 + 0.135 = 4.465 s, about 4.5 s. ENGAGE: why design for the 90th percentile and not the mean? (Expected: a speed chosen for the mean is overrun by every condition slower than the mean.) WARN: these are means of conditions; an individual operator's tail is longer still, and the site must measure its own operators. FLEX: [CORE].}
\end{frame}

\begin{frame}{No brake buys back a slow human: 0.836 m/s at most}
  \begin{columns}[T]
    \begin{column}{0.38\textwidth}
      \begin{mlsyscard}{mlsysaccent}
        \small\textbf{Given (running machine)}\\[0.05cm]
        Room to brake $D_{\text{clear}}-\delta = 0.91$~m\\
        Deceleration $a_{\text{brake}} = 2$~m/s$^2$\\
        Jerk limit $j_{\max} = 7.11$~m/s$^3$\\
        Machine delay $\tau_{\text{delay}} = 133.6$~ms\\
        Takeover $T_{\Sigma,1} = 605$~ms
      \end{mlsyscard}
    \end{column}
    \begin{column}{0.60\textwidth}
      \small
      \wstep{2}{Time at full speed: $b = 0.1336+0.605+\tfrac{2}{2\times7.11} = 0.8792$~s}
      \wstep{3}{$v_{\max} = a_{\text{brake}}\bigl[\sqrt{b^2+2(0.91)/a_{\text{brake}}}-b\bigr] = 2\,(1.2973-0.8792)$}
      \wstep{3}{$v_{\max} = \mathbf{0.836}$~m/s, below the 1.3~m/s aisle speed}
      \wstep{4}{Perfect brake: $0.91/(0.1336+0.605) = \mathbf{1.23}$~m/s, still below 1.3}
      \wstep{5}{With the 63.3~ms admission allocation: \textbf{\textcolor{mlsysaccent}{0.797~m/s}}}
    \end{column}
  \end{columns}
  \mlsyscredit{Supervisory Intervention, Napkin Math 14.2 and eq.~3 (illustrative and derived values); dispatch-to-admission allocation from Lecture 13.}
  \note{[2.5 min] LINK: Lecture 2 owns the velocity ceiling as the root of the stopping inequality; this adds the human's delay to it. NARRATE: Step 2: every interval at full speed adds up: the machine's 133.6 ms, the human's 605 ms, and half the jerk ramp, a over 2j, 140.6 ms, so b = 0.8792 s. Step 3: the quadratic root gives 0.836 m/s. Check the branch: the root is above a squared over j, 0.562 m/s, so the quadratic form applies. Step 4: remove braking limits entirely and the ceiling is the room divided by the delays, 1.23 m/s, still under 1.3. Step 5: charging Lecture 13's 63.3 ms lowers it to 0.797 m/s. All derived. ENGAGE: what are the two ways to make an in-loop takeover safe? (Expected: slow down before asking, or stop before asking.) WARN: a designer who stops at 0.84 m/s and picks 0.8 misses the complete budget's deadline by 5.0 ms once the 63.3 ms is charged. FLEX: [OPTIONAL]. If cut, state 0.836 and 1.23 m/s on the next slide.}
\end{frame}

\begin{frame}{The deadline is the same budget solved for the wait}
  \centering
  \mlsyssteps[width=0.88\textwidth,height=4.6cm,keepaspectratio]{deadline_step}{3}\\[0.05cm]
  {\small $T_{\text{deadline,fallback}} = \dfrac{D_{\text{clear}}-\delta-d_{\text{suffix}}(v)}{v}-\tau_{\text{delay}}$\quad Configured timeouts are clipped to it.}
  \mlsyscredit{Supervisory Intervention, eq.~4 and Napkin Math 14.3 (derived values of the running machine).}
  \note{[3 min] NARRATE: fix the speed and solve the stopping inequality for the wait. Step 1: the green curve is the longest the machine may wait after a request before its own stop must begin. At 1.3 m/s it is 100.8 ms, and 37.5 ms after the admission allocation. Step 2: the in-loop takeover needs 605 ms. At 0.75 m/s the deadline is 751.6 ms; the takeover leaves 146.6 ms, and 83.3 ms after charging 63.3 ms. The design slows to 0.75 m/s, the 0.797 m/s ceiling rounded down to a 0.05 m/s step, before an anticipated in-loop request. Slowing multiplies the allowance 7.5 times, because both the stop and the travel per millisecond shrink. Step 3: the out-of-loop 4.5 s needs a crawl; at 0.15 m/s the deadline is 5.79 s and leaves 1.32 s. A 99th-percentile operator (7.1 s) would still overrun, so the robust design stops first and hands over from standstill, where waiting costs time but no clearance. All derived. ENGAGE: road rules let an automated lane-keeping car wait 10 s (UN Regulation 157). Why can the robot not? (Expected: 10 s at 1.3 m/s is 13 m.) WARN: a hidden pallet leaves no time for a planned slow-down; the machine answers it with its own stop and asks from standstill. FLEX: [CORE].}
\end{frame}

\begin{frame}{Estimate: how fast may a braking driver be asked?}
  \mlsysactivity{Estimate}{3 min}
  \begin{columns}[c]
    \begin{column}{0.56\textwidth}
      \begin{mlsyscard}{mlsysaccent}
        \footnotesize\textbf{Given}\\[0.03cm]
        Expected brake signal: 750~ms (Green 2000), plus the 135~ms console term.\\
        $b = \tau_{\text{delay}} + T_\Sigma + a/(2j)$\\ $v_{\max} = a\bigl[\sqrt{b^2+0.91}-b\bigr]$
      \end{mlsyscard}
      \only<1|handout:0>{\normalsize In pairs: the request speed?}%
      \uncover<2->{\small $T_\Sigma = 885$~ms, $b = 1.1592$~s\\[0.05cm]
      $v_{\max} = \mathbf{0.68}$~m/s, as in table 2}
    \end{column}
    \begin{column}{0.41\textwidth}
      \centering\safeimg[width=\textwidth,height=2.6cm,keepaspectratio]{uber-atg-volvo-xc90.jpg}\\[0.03cm]
      {\footnotesize An Uber self-driving test car, the Volvo XC90 model of the Tempe crash, 2017}
    \end{column}
  \end{columns}
  \mlsyscredit{Supervisory Intervention, Checkpoint 14.1 and table 2 (derived); Green 2000. Photo: zombieite, CC BY 2.0, Wikimedia Commons.}
  \note{[3 min] ENGAGE: two minutes in pairs. Step 2: the wait is 750 + 135 = 885 ms; b = 1.1592 s; the root is 0.684 m/s, the table's 0.68 m/s (derived). A driver who expects the brake signal is the in-loop case at its best; a surprise raises the response to about 1.5 s (Green 2000) and the ceiling to 0.45 m/s. NARRATE: the photo shows the same vehicle model and program as the 2018 Tempe crash, which Lecture 1 owns. There the safety driver was looking at a phone, and the system had disabled automatic emergency braking and relied on the person to intervene (NTSB HAR-19/03). For today, an operator-monitoring contract must be backed by a local stop that works whether or not the person is watching. WARN: this photo is the same model, not the crash vehicle. FLEX: [CORE].}
\end{frame}

% =============================================================================
% SECTION: HANDOVER
% =============================================================================
\mlsysnote{[0.5 min] [CORE] Now the mechanism: how authority is bounded in time and moved between sources without two writers on one actuator.}
\mlsysdivider{How does authority change hands without two writers?}{Supervisory Intervention, \S14.3}
\section{Handover}

\begin{frame}{A three-position grip stops on release and on a squeeze}
  \begin{columns}[c]
    \begin{column}{0.32\textwidth}
      \centering\safeimg[width=\textwidth,height=4.2cm,keepaspectratio]{abb-teach-pendant.jpg}\\[0.03cm]
      {\footnotesize An ABB teach pendant; its rear grip proves physical presence}
    \end{column}
    \begin{column}{0.65\textwidth}
      \footnotesize
      \begin{tabular}{@{}ll@{}}\toprule
        \textbf{Device or rule} & \textbf{Behavior} \\\midrule
        Grip released (position 1) & stop requested \\
        Grip in the middle (position 2) & motion enabled \\
        Grip squeezed hard (position 3) & stop requested \\
        Manual speed, ISO 10218-1:2025 & at most 250~mm/s \\
        Emergency stop, ISO 13850 & latches until local reset \\\bottomrule
      \end{tabular}\\[0.2cm]
      {\small At 250~mm/s, a 470~ms human response is \textbf{117.5~mm} of travel. The remote takeover at 1.3~m/s is \textbf{786.5~mm}.}
    \end{column}
  \end{columns}
  \mlsyscredit{Supervisory Intervention, \S14.3.1 and footnote 5 (ISO 10218-1:2025, ISO 13849-1, ISO 13850). Photo: Auledas, CC BY-SA 4.0, Wikimedia Commons.}
  \note{[2 min] NARRATE: a technician beside a robot jogs it from a teach pendant like this one. Motion needs a hand on the enabling grip on the back, and the grip has three positions. A startled person either lets go or squeezes hard; a two-position switch would miss the squeeze, so only the middle position enables motion, and returning to the middle after a squeeze must not restart motion. The standards cap the tool speed in this mode at 250 mm/s, so a 470 ms human response costs 117.5 mm (derived), against 786.5 mm for the remote takeover at aisle speed. The red button latches and only a deliberate local reset clears it, which makes a human stop sticky in a way a network command is not. ENGAGE: why does the remote console need no grip? (Expected: authenticated heartbeats keep its lease alive; a site may add a grip as a second condition.) WARN: releasing the grip only requests the stop the safety circuit selects; Safe Torque Off alone does not hold a loaded arm. FLEX: [CORE].}
\end{frame}

\begin{frame}[fragile]{Authority leases lapse on the MCU's own clock}
  \begin{columns}[T]
    \begin{column}{0.50\textwidth}
      {\footnotesize\textbf{ROS 2 teleop\_twist\_joy, on enable-button release}}\par\vspace{0.05cm}
      {\fontsize{6.8}{8.2}\selectfont
\begin{semiverbatim}
\} else \{
  // When enable button is released, immediately
  // send a single no-motion command in order to
  // stop the robot.
  if (!sent_disable_msg) \{
    auto cmd_vel_msg =
      std::make_unique<geometry_msgs::msg::Twist>();
    cmd_vel_pub->publish(std::move(cmd_vel_msg));
    sent_disable_msg = true;
  \}
\}
\end{semiverbatim}}
      {\small A lost zero command never arrives. The lease expires anyway; a grip adds physical presence.}
    \end{column}
    \begin{column}{0.48\textwidth}
      \centering\safeimg[width=\textwidth,height=4.5cm,keepaspectratio]{leases.pdf}
    \end{column}
  \end{columns}
  \mlsyscredit{github.com/ros2/teleop\_twist\_joy, src/teleop\_twist\_joy.cpp at 37a30e7 (abridged; comment rewrapped). Supervisory Intervention, \S14.3.1, \S14.3.4, eq.~5.}
  \note{[2.5 min] LINK: Lecture 4 gave the chunk lease and authority-expiry timers to the safety microcontroller; Lecture 3's chunk policy renews every 50 ms. NARRATE: left is real code from the standard ROS 2 joystick teleoperation node. When the operator releases the enable button, it publishes one zero-velocity message, once. If that message is dropped, nothing on this side sends another. Right, the chapter's design: each proposer holds a lease on the safety microcontroller's own clock. The policy's chunk lease is 60 ms; the operator's manual lease is 100 ms, renewed by authenticated heartbeats. When renewals stop, the clock runs out with no packet needed, and the arbiter picks the state-matched HOLD or STOP rung (Lecture 12's fallback ladder). Return to motion takes standstill and a fresh handshake; resumed packets alone do not restore it. ENGAGE: what bounds the manual lease from below and from above? (Expected: below, the heartbeat period plus jitter and one lost packet; above, its extra 40 ms over the chunk lease is travel, 30 mm at 0.75 m/s.) WARN: a heartbeat renews authority, not evidence; an old observation stays old. FLEX: [CORE].}
\end{frame}

\begin{frame}{Each channel is in exactly one authority state per tick}
  \centering
  \mlsyssteps[width=0.95\textwidth,height=5.1cm,keepaspectratio]{fsm_step}{3}\\[0.02cm]
  {\small The enforcer is the only writer to the drives in every state.}
  \mlsyscredit{Supervisory Intervention, \S14.3.2, table 4 and figure 5 (redrawn).}
  \note{[3 min] NARRATE: Step 1: the arbiter resolves each channel to one state every 1 ms tick. In Autonomous the policy holds the proposal token. A request moves it to Handover-Pending, where the policy keeps proposing under the clearance-derived timer. Commit moves it to Blend, where no proposer holds the token and the enforcer drives a smooth profile between latched endpoints. When the profile completes under a valid manual lease, the operator receives the token in Human-Manual. Step 2: every active state has an edge to Fallback, on lease lapse, timeout, or invalid state, so no failure of the network or the person leaves a channel without an owner. The only way out of Fallback is a reset from rest. Step 3: Protective sits outside the protocol; a hardwired protective input preempts every state. ENGAGE: why does Fallback not return to the policy that asked for help? (Expected: that policy asked because it was failing; the validated stop rung is the safe owner.) WARN: CAN-bus arbitration serializes frames but does not decide authority; two valid streams to one drive make it chatter. FLEX: [CORE].}
\end{frame}

\begin{frame}{Four phases commit ownership on a single 1 ms tick}
  \centering
  \safeimg[width=0.92\textwidth,height=5.0cm,keepaspectratio]{handshake.pdf}\\[0.05cm]
  {\small The handover handshake keeps a single writer. Chosen budgets: 10 + 5 + 1 + 15 = \textbf{31~ms}.}
  \mlsyscredit{Supervisory Intervention, \S14.3.3, table 5 (chosen budgets).}
  \note{[2 min] NARRATE: Request carries an authenticated role, a nonce, and a fresh timestamp, and starts the clearance-derived countdown. Acknowledge checks fresh state, a feasible manual target, and full stopping reserve. Commit happens on the next tick: the policy's token is revoked on the named channels, both endpoints are latched, and the enforcer becomes the sole writer. Confirm tells the console; if that message is lost, the console alarms, but the commit on the robot stands. ENGAGE: why separate the local commit from the network confirmation? (Expected: so a dropped packet leaves only the console uncertain and never the drives.) WARN: the 31 ms is a chosen budget, not a measurement, and it is part of the 605 ms, not added to it. FLEX: [OPTIONAL]. If cut, show the 31 ms in the timing strip.}
\end{frame}

\begin{frame}{Vote: what sets the handover timeout here?}
  \mlsysactivity{Concept question}{2.5 min}
  \begin{mlsyscard}{mlsysaccent}
    \small The policy asks for help at 1.3~m/s. The configured protocol ceiling is 500~ms. The clearance deadline, after the 63.3~ms admission allocation, is 37.5~ms. The arbiter runs a 1~ms tick.
  \end{mlsyscard}
  \vspace{0.2cm}
  {\normalsize\choice{A}{Wait 500 ms for the operator}\\[0.12cm]
  \answerat{2}{B}{Fall back after 37 whole ticks}\\[0.12cm]
  \choice{C}{Wait until the console confirms}}\\[0.2cm]
  \uncover<2->{\small $T_{\text{timeout}} = \max(0, \min(500, 37.5))$~ms. The clearance deadline is tighter by \textbf{462.5~ms}.}
  \mlsyscredit{Supervisory Intervention, eq.~4 and Checkpoint 14.2 (derived values).}
  \note{[2.5 min] ENGAGE: vote, discuss with a neighbor for 30 seconds, vote again. Step 2: B. The timeout is the smaller of the configured ceiling and the clearance deadline, so 37.5 ms; the arbiter dispatches fallback by the last whole tick, 37 ticks (derived). The 500 ms ceiling is only an outer protocol limit. C waits on the network, which the protocol never does. WARN: students treat a timeout as a setting. It is a function of the current speed and clearance: the same 500 ms window waited out at 1.3 m/s gives a stop of 1.62 m, 0.52 m past the rack end, while at a 0.15 m/s crawl the whole stop is 0.31 m (derived, the chapter's pitfall). Never extend a running countdown on an optimistic recalculation; only tighten it. FLEX: [CORE].}
\end{frame}

% =============================================================================
% SECTION: TRANSFER
% =============================================================================
\mlsysnote{[0.5 min] [CORE] Ownership now changes on one tick. The command itself must not.}
\mlsysdivider{Why can the command not switch in one tick?}{Supervisory Intervention, \S14.4 and \S14.5}
\section{Transfer}

\begin{frame}{A one-tick step asks the drivetrain for 2,000 m/s\texorpdfstring{$^3$}{3}}
  \begin{columns}[c]
    \begin{column}{0.66\textwidth}
      \centering\mlsyssteps[width=\textwidth,height=5.0cm,keepaspectratio]{quintic_step}{2}
    \end{column}
    \begin{column}{0.31\textwidth}
      \centering\safeimg[width=\textwidth,height=2.6cm,keepaspectratio]{strain-wave-gear.jpg}\\[0.03cm]
      {\footnotesize A strain wave gear set, the reducer in many arm joints}\\[0.2cm]
      {\small Bumpless transfer: $\alpha(s)=10s^3-15s^4+6s^5$, peak slope 1.875.}
    \end{column}
  \end{columns}
  \mlsyscredit{Supervisory Intervention, \S14.4.1 and \S14.4.2, eq.~6 and 7 (derived). Photo: Pieceofmetalwork, CC BY-SA 4.0, Wikimedia Commons.}
  \note{[2.5 min] LINK: Lecture 11 joined trajectories with C2 seams, continuous in value, rate, and acceleration; the enforcer reuses that construction at commit. NARRATE: Step 1: at commit the enforcer has two latched commands, the policy's last admitted one and the operator's checked target. Switching between them in one 1 ms tick turns a 2 m/s squared change into a jerk of 2,000 m/s cubed for that tick. A torque step snaps gears across their backlash, rings the drivetrain, and sags the shared supply rail. Step 2: the quintic smoothstep spreads the change over a window so the command and its first two derivatives are continuous; its slope peaks at 1.875 over the window, so the window must be at least 1.875 times the change over the rate limit. At the base's 7.11 m/s cubed that is 527.34 ms (derived). ENGAGE: what does a longer window buy and cost? (Expected: lower peak jerk; later arrival of the operator's command and more travel.) WARN: a smooth command is not a measured smooth response; motor current, encoders, and the inertial sensor must confirm it. FLEX: [CORE].}
\end{frame}

\begin{frame}{Estimate: how long must the arm's blend last?}
  \mlsysactivity{Estimate}{2.5 min}
  \begin{columns}[T]
    \begin{column}{0.42\textwidth}
      \begin{mlsyscard}{mlsysaccent}
        \small\textbf{Given (illustrative)}\\[0.05cm]
        Arm joint torque change 15~N$\cdot$m\\
        Torque-rate ceiling 750~N$\cdot$m/s\\
        $\tau_{\text{blend}} \ge 1.875\,|\Delta u| / \text{rate ceiling}$\\
        Permission tick 1~ms
      \end{mlsyscard}
    \end{column}
    \begin{column}{0.55\textwidth}
      \only<1|handout:0>{\normalsize In pairs: the shortest arm blend, in whole ticks?}%
      \uncover<2->{\small $1.875 \times 15 / 750 = 37.5$~ms, so \textbf{38} ticks.\\[0.15cm]
      Base: $1.875 \times 2 / 7.11 = 527.34$~ms, so \textbf{528} ticks, at any speed.\\[0.15cm]
      From 0.75~m/s, a smoothed stop needs 328.4~mm; the jerk-limited stop 246.1~mm; an ideal step 140.6~mm.}
    \end{column}
  \end{columns}
  \mlsyscredit{Supervisory Intervention, \S14.4.2 and table 6 (derived values).}
  \note{[2.5 min] ENGAGE: one minute. Step 2: 37.5 ms, rounded up to 38 ticks; the base window is 527.34 ms, 528 ticks, and it does not depend on speed because the jerk bound belongs to the loaded base. NARRATE: smoothing costs clearance. From the 0.75 m/s request speed, a blend into braking followed by full braking stops in 328.4 mm, against 246.1 mm for Lecture 11's jerk-limited resident stop and 140.6 mm for an ideal step (derived). For a braking command the resident stop already is the bumpless transfer, so the base takes it, and the window governs only non-stopping changes, such as the first manual command after a stop. WARN: an unblended emergency stop is shorter only on paper; wheel lockup or a thrown tote can lengthen the real stop. FLEX: [CORE].}
\end{frame}

\begin{frame}{The in-loop session at 0.75 m/s closes by 32.5 mm}
  \centering
  \safeimg[width=0.88\textwidth,height=4.5cm,keepaspectratio]{session.pdf}\\[0.05cm]
  {\small The base takes its resident stop. Only the arm blends.}
  \mlsyscredit{Supervisory Intervention, \S14.4.6, table 7 (derived; the 157 mm blend reserve is the chapter's value).}
  \note{[2 min] NARRATE: every interval priced so far now draws on one reserve. At 0.75 m/s the finished budget leaves 563.7 mm. The 605 ms takeover, handshake included, travels 453.8 mm, leaving 110.0 mm. The 63.3 ms admission allocation takes 47.5 mm, leaving 62.5 mm. Once the operator holds the channel, the 100 ms manual lease replaces the 60 ms chunk lease in the delay, and its extra 40 ms costs 30 mm, leaving 32.5 mm. A braking blend on the base would need up to 157 mm, so it does not fit. All derived except the 157 mm, which is the chapter's. ENGAGE: which configurations are worth qualifying instead? (Expected: the 0.15 m/s crawl, which keeps 198.4 mm after a 4.5 s takeover, or standstill.) WARN: the live check uses the current clearance and only the unelapsed parts of each interval; it never charges the elapsed delay twice. Remote assistance is never a safety layer. LINK: how many machines one assistant can serve is Lecture 17's question. FLEX: [OPTIONAL].}
\end{frame}

\begin{frame}{Averaging two good commands can hit the pallet}
  \begin{columns}[c]
    \begin{column}{0.50\textwidth}
      \centering\safeimg[width=\textwidth,height=5.2cm,keepaspectratio]{pallet.pdf}
    \end{column}
    \begin{column}{0.47\textwidth}
      \centering\safeimg[width=\textwidth,height=3.2cm,keepaspectratio]{da-vinci-console.jpg}\\[0.03cm]
      {\footnotesize A da Vinci surgeon console: hand controllers drive follower arms through a computer}\\[0.2cm]
      {\small Blending works when the task has one goal and every blend is admissible.}
    \end{column}
  \end{columns}
  \mlsyscredit{Supervisory Intervention, \S14.3 (Dragan and Srinivasa 2013). Photo: Cmglee, CC BY-SA 3.0, Wikimedia Commons.}
  \note{[1.5 min] NARRATE: shared autonomy keeps both proposers active and blends them. The robot infers the operator's goal and mixes assistance with manual control; user studies show it works while the robot predicts that goal with high confidence (Dragan and Srinivasa 2013). Around an obstacle there are two good paths. If the operator goes left and the policy goes right, their average goes straight into the pallet, a third command nobody proposed and nobody checked. The surgical console on the right is a leader console of this kind, used for fine single-goal motion. ENGAGE: name a task where blending is safe. (Expected: fine adjustment of a gripper pose toward one target.) WARN: the chapter selects one source per tick; the enforcer still filters whatever it selects. FLEX: [OPTIONAL].}
\end{frame}

\begin{frame}{Diagnose: can valid cryptography delay a stop?}
  \mlsysactivity{Diagnose}{2.5 min}
  \begin{columns}[c]
    \begin{column}{0.58\textwidth}
      \centering\safeimg[width=\textwidth,height=4.4cm,keepaspectratio]{flood.pdf}
    \end{column}
    \begin{column}{0.40\textwidth}
      \small Motion commands carry a digital signature (Ed25519). A faulty radio floods the MCU with forged requests, and a genuine veto waits in the same queue.\\[0.15cm]
      \only<1|handout:0>{\normalsize What broke, and what fixes it?}%
      \uncover<2->{\textbf{55} forged packets miss the 100.8~ms deadline; \textbf{20} miss 37.5~ms. Fix: a rate limiter before the verifier; hardwired stops skip it.}
    \end{column}
  \end{columns}
  \mlsyscredit{Supervisory Intervention, \S14.5 and \S14.9; 1.8 ms per Ed25519 check is the chapter's estimate from published Cortex-M4 cycle counts (Owens et al. 2024).}
  \note{[2.5 min] NARRATE: a motion command injects energy, so it must be authenticated, attributable, and fresh; a protective stop never waits for verification. Ed25519 verification is estimated at 1.8 ms per packet on the safety microcontroller, from published Cortex-M4 cycle counts scaled to its clock; one check costs 2.34 mm of travel at aisle speed. ENGAGE: give one minute. Step 2: the verifier spends 1.8 ms on each well-formed forgery before it reaches the veto. With 55 queued ahead, the veto is verified at 100.8 ms, past the 100.78 ms deadline; with the admission allocation charged, 20 suffice. Sustaining it takes 546 packets per second, trivial for any radio (derived). The fault is availability, not secrecy. The fix is a per-source token-bucket limiter ahead of the verifier, sized from the stopping budget, with emergency stops and lease timers on hardwired paths that never share the queue. WARN: a 100 ms freshness window is 130 mm at 1.3 m/s and leaves 1.0 mm of the spare, so freshness is a distance budget too. FLEX: [CORE].}
\end{frame}

% =============================================================================
% SECTION: RECORD
% =============================================================================
\mlsysnote{[0.5 min] [CORE] The bus now has one writer per channel. The person's picture of who drives may still be wrong, and the record must say what happened.}
\mlsysdivider{Who held each channel, and what may be learned?}{Supervisory Intervention, \S14.6 to \S14.8}
\section{Record}

\begin{frame}{Pilots act on the mode they believe the autopilot holds}
  \begin{columns}[c]
    \begin{column}{0.56\textwidth}
      \centering\safeimg[width=\textwidth,height=4.4cm,keepaspectratio]{a320-primary-flight-display.jpg}\\[0.03cm]
      {\footnotesize An A320 flight display in cruise. Top row: the active mode of each channel, and which automation is engaged.}
    \end{column}
    \begin{column}{0.41\textwidth}
      \small\textbf{Mode confusion:} the mode a supervisor believes differs from the actual one, so a correct action for the believed mode is hazardous in the real one.\\[0.2cm]
      {\footnotesize Sarter and Woods, 1995, from cockpit automation.}
    \end{column}
  \end{columns}
  \mlsyscredit{Supervisory Intervention, \S14.6 and Definition 14.3. Photo: Olivier Cleynen, CC0, Wikimedia Commons.}
  \note{[1.5 min] NARRATE: aviation met this failure first. Sarter and Woods traced accidents in automated cockpits to pilots whose inputs were right for the mode they assumed and wrong for the mode the aircraft was in. Airliners answer with a flight mode annunciator: the top row of this display names the active speed, vertical, and lateral modes and the engaged autopilot and autothrust. Every channel is announced all the time. ENGAGE: is mode confusion an arbitration fault? (Expected: no; arbitration is solved by one writer on the bus, while mode confusion lives in the person's model of the machine.) WARN: tracking where an operator looks says nothing about whether they know which channel the policy still drives. FLEX: [OPTIONAL].}
\end{frame}

\begin{frame}{The base kept driving while the operator aligned the arm}
  \centering
  \mlsyssteps[width=0.95\textwidth,height=5.5cm,keepaspectratio]{channels_step}{2}
  \mlsyscredit{Supervisory Intervention, \S14.6 (the wet shrink-wrap case; derived travel at the aisle speed).}
  \note{[2.5 min] NARRATE: wet shrink-wrap confuses the depth camera, and the runtime asks a remote assistant to reposition the gripper. Step 1: the console grants the arm channel only. The assistant, focused on the fingers, assumes the base is parked, but the base channel was never transferred and the policy keeps driving at 1.3 m/s. Over the 605 ms window the base covers 786.5 mm, 6.0 times the 131 mm spare (derived). Every component met its specification. Step 2: two rules. Every Request, Acknowledge, and Commit names its channel mask, and the console announces the owner of every channel. A channel left out of the transfer may keep its autonomous owner only while the console says so and the operator acknowledged that state; otherwise the commit puts it in a certified hold with the brake on. ENGAGE: which rule would have stopped this collision? (Expected: either one; the second is the safe default.) WARN: one machine-wide mode flag cannot show a split handover. FLEX: [CORE].}
\end{frame}

\begin{frame}{Flight 2904: an interlock refused braking for 9 s}
  \begin{columns}[c]
    \begin{column}{0.67\textwidth}
      \centering\safeimg[width=\textwidth,height=4.2cm,keepaspectratio]{lh2904.pdf}
    \end{column}
    \begin{column}{0.30\textwidth}
      \centering\safeimg[width=\textwidth,height=2.3cm,keepaspectratio]{lufthansa-a320-d-aipn.jpg}\\[0.03cm]
      {\footnotesize D-AIPN, the A320 of flight 2904, at Heathrow in 1991}\\[0.15cm]
      {\small Warsaw, 1993: overran at 72 knots; two died.}
    \end{column}
  \end{columns}
  \mlsyscredit{Supervisory Intervention, War Story 14.1, from the Warsaw accident report (1994). Photo: Dick Gilbert, CC BY 2.0, Wikimedia Commons.}
  \note{[2.5 min] NARRATE: Lufthansa 2904 landed fast, far down a waterlogged runway with an unreported tailwind. The crew commanded brakes and reversers at once. The ground spoilers and reversers required both main struts compressed or the wheels spun up; the left strut stayed uncompressed for 9 s and the wheels aquaplaned, so the logic decided the aircraft was still flying. Spoilers and reversers stayed locked out until 9 s and wheel braking until 13 s. The aircraft overran at 72 knots, struck an embankment, and burned; two people died. The report names the interlock design alongside crew actions and does not claim an override alone would have prevented the overrun. ENGAGE: the robot version? (Expected: a wheel-slip observer or a docking-latch check that gates a supervisor's braking command.) WARN: the lesson is not to bypass interlocks. Annunciate the unsatisfied predicate, name it, and provide a certified fallback deceleration path for sensor disagreement. FLEX: [CORE].}
\end{frame}

\begin{frame}{True or false: takeover should bypass the filter?}
  \mlsysactivity{Concept question}{2 min}
  \centering
  {\normalsize ``A person who takes over sees what the model missed, so manual commands should go straight to the drives.''}\\[0.4cm]
  {\large\choice{A}{True}\qquad\answerat{2}{B}{False}}\\[0.4cm]
  \uncover<2->{\small The operator gains the right to propose. Momentum, friction, and the rack do not change, and the operator's response is the slowest term in the loop.}
  \mlsyscredit{Supervisory Intervention, \S14.9 (fallacy 1).}
  \note{[2 min] ENGAGE: vote, then ask one student from each side to argue. Step 2: false. A takeover changes only who proposes. The enforcer projects a manual command onto the admissible set, or rejects it and falls back, exactly as for the policy. The manual authority also lives under a 100 ms lease, so a dropped link cannot leave a stale velocity running. Flight 2904 seems to argue the other way, but the chapter draws the opposite lesson: annunciate why the machine refuses, and provide a certified fallback stop, rather than an unchecked bypass. WARN: an operator who cannot see why the robot ignores the joystick pushes harder; the console must show the active constraint. FLEX: [CORE].}
\end{frame}

\begin{frame}{The authority record logs who held each channel every tick}
  \begin{columns}[T]
    \begin{column}{0.58\textwidth}
      {\footnotesize\textbf{Per-tick evidence tuple, 1 kHz}}\par\vspace{0.06cm}
      \begin{logblock}[7.2]
        \logrow{t\_mcu}{MCU hardware clock, ns}
        \logrow{authority\_state}{AUTO ... PROTECTIVE}
        \logrow{measured\_state}{joints, currents, inertia}
        \logrow{proposed\_command}{policy or operator}
        \logrow{admitted\_command}{after the enforcer}
        \logrow{dispatched\_torque}{to the drives}
        \logrow{operator\_id}{authenticated holder}
        \logrow{lease\_remaining}{chunk or manual}
      \end{logblock}
      \vspace{0.1cm}
      {\footnotesize Each entry hashes the one before it; a missing tick is a declared gap.}
    \end{column}
    \begin{column}{0.39\textwidth}
      \centering\safeimg[width=\textwidth,height=2.4cm,keepaspectratio]{flight-data-recorder.jpg}\\[0.03cm]
      {\footnotesize A flight data recorder (NTSB photo)}\\[0.15cm]
      {\small A car's event recorder keeps \textbf{10} samples of the last 5~s (2~Hz); the authority record keeps \textbf{5,000}.}
    \end{column}
  \end{columns}
  \mlsyscredit{Supervisory Intervention, \S14.7 and \S14.8 (tuple abridged); 49 CFR Part 563 (2024). Photo: NTSB, public domain, Wikimedia Commons.}
  \note{[2.5 min] LINK: Lecture 1 demanded synchronized records on both sides of the causal boundary; Lecture 12's enforcement record supplies the fallback rungs and Lecture 13's placement record the timing bounds. NARRATE: this lecture's record, the twelfth in the record chain. A sealed static manifest, loaded before the drives energize, binds each channel to its allowed roles, preemption tier, lease, rate ceiling, and blend window. Every tick then logs this tuple. Without the proposal, a torque cannot show who asked for it; without the admitted command, the log cannot show whether the barrier changed it; without the lease remaining, it cannot show whether authority had already lapsed. Entries are hash-chained, and a missing tick is a declared gap. A passenger car's event data recorder keeps 5 s at 2 Hz, 10 samples, rising to 20 s at 10 Hz from 2027; the 5 s pre-trigger window here holds 5,000 tuples (derived). ENGAGE: why not reuse the car standard? (Expected: one sample per 100 ms is 100 permission ticks, too coarse to order a takeover against a contact.) WARN: the record shows who held each channel, not who decided the machine should be there; that approval record is Lecture 17's. The manifest's four rules, A1 to A4, are the fault cases Lecture 15 injects. FLEX: [CORE].}
\end{frame}

\begin{frame}{Estimate: how many logged ticks may train?}
  \mlsysactivity{Estimate}{3 min}
  \centering
  \mlsyssteps[width=0.92\textwidth,height=4.0cm,keepaspectratio]{masks_step}{2}\\[0.02cm]
  {\small With 4.5~s of drift, a 38~ms blend, 1~ms ticks, and no automated expert, how many ticks may train?}\\[0.05cm]
  \uncover<2->{\small Drift \textbf{4,500} ticks and blend \textbf{38} ticks, of which \textbf{0} may train.}
  \mlsyscredit{Supervisory Intervention, \S14.8.3, eq.~8 and table 8; Sirius: Liu et al., RSS 2023.}
  \note{[3 min] LINK: Lecture 5 owns the curation rules for intervention data; the record supplies the timestamps those rules key on. ENGAGE: ask for the counts. Step 1 answer: interval 1, the drift, is 4,500 ticks; interval 2, the blend, is 38 ticks (derived). Without an automated expert, none of them may train an imitation policy: drift commands are the failure, and the blend is the enforcer's spline, which belongs to neither the policy nor the person. Interval 3, the operator hunting while re-orienting, is masked until tracking error and command rate stay inside limits for a dwell; drivers resuming control have needed 35 to 40 s before lateral steering settled (Merat et al. 2014). Only interval 4 trains as a demonstration. Step 2: Sirius zeroes a fixed 2 s window before the intervention. For a 4.5 s takeover that leaves 2,500 drift ticks in as if they were good (derived). WARN: training on raw rescue logs teaches the policy the drift and the startled corrections that made the takeover necessary. Other methods (HG-DAgger, RLIF, HIL-SERL, RECAP) keep some intervals, but all need the owner of every sample. FLEX: [CORE].}
\end{frame}

% =============================================================================
% CLOSE
% =============================================================================
\begin{frame}{Retrieval: what does the takeover count leave out?}
  \mlsysactivity{Retrieval}{2 min}
  \begin{columns}[c]
    \begin{column}{0.52\textwidth}
      \centering\safeimg[width=\textwidth,height=3.6cm,keepaspectratio]{dmv_step2.pdf}
    \end{column}
    \begin{column}{0.46\textwidth}
      \only<1|handout:0>{\normalsize Close your notes. Name the four things a takeover needs that the filing does not record.}%
      \uncover<2->{\small 1.\enspace A deadline: 605~ms in the loop caps the request at 0.84~m/s\\
      2.\enspace One writer per channel: leases, six states, a 31~ms handshake\\
      3.\enspace A smooth transfer: 528~ms base, 38~ms arm\\
      4.\enspace A record: 5,000 tuples in 5~s, and masks}
    \end{column}
  \end{columns}
  \mlsyscredit{California DMV disengagement reports (Waymo); Supervisory Intervention, \S14.10.}
  \note{[2 min] Back to the hook. ENGAGE: 60 seconds of writing. Step 2: the four answers. A deadline: a takeover is travel before braking; the in-loop response caps the request speed at 0.84 m/s and the out-of-loop one at 0.19 m/s, so the machine slows, crawls, or stops before it asks. One writer: leases on the MCU clock, one authority state per channel per tick, and a four-phase handshake that commits on one tick. A smooth transfer: the quintic blend bounds the command rate and costs clearance. A record: the per-tick tuple says who held each channel, and its timestamps decide which samples may train. The fleet in the plot still takes 174 takeovers a year; human authority is permanent, so each of those four is part of the design. FLEX: [CORE].}
\end{frame}

\begin{frame}{Takeaways}
  \thumbline{thumb_strip_step3.pdf}{A takeover is travel before braking: 605~ms is 786.5~mm at 1.3~m/s, so the machine slows, crawls, or stops before it asks.}\\[0.15cm]
  \thumbline{thumb_fsm_step3.pdf}{Leases on the MCU clock end silent authority; each channel has one state and one writer every tick.}\\[0.15cm]
  \thumbline{thumb_quintic_step2.pdf}{A quintic blend replaces a 2,000~m/s$^3$ step, and its window costs clearance.}\\[0.15cm]
  \thumbline{thumb_masks_step1.pdf}{Announce every channel's owner, and log it per tick; only stable recovery trains.}
  \note{[1 min] [CORE] One line per segment. Ask students which slide taught each line. If short on time, read the first and the second.}
\end{frame}

\begin{frame}{Open questions}
  \begin{columns}[c]
    \begin{column}{0.62\textwidth}
      \small
      \begin{itemize}\setlength{\itemsep}{0.12cm}
        \item If takeovers become rare, operators lose practice exactly when a takeover is hardest. How should a fleet keep them ready?
        \item Should a robot ever accept a moving takeover, or always stop first?
        \item Who should own the record of a takeover: the operator, the fleet company, or the site?
      \end{itemize}
    \end{column}
    \begin{column}{0.36\textwidth}
      \centering\safeimg[width=\textwidth,height=3.2cm,keepaspectratio]{iss-robotics-workstation.jpg}\\[0.03cm]
      {\footnotesize Astronauts train for arm operations they perform a few times a year}
    \end{column}
  \end{columns}
  \mlsyscredit{Photo: NASA, public domain, Wikimedia Commons.}
  \note{[1 min] [OPTIONAL] No single answer. The first is Bainbridge's irony of automation (1983): automation leaves the person the rare, hardest cases with skills it has let fade. The second weighs availability against the out-of-loop tail of 4.5 to 7 s. The third anticipates Lecture 16's release argument and Lecture 17's approval record.}
\end{frame}

\begin{frame}{Read and explore}
  \small
  \begin{itemize}\setlength{\itemsep}{0.02cm}
    \item \textbf{Read:} Supervisory Intervention, \S14.1 to \S14.10; Checkpoints 14.1 to 14.3; Napkin Math 14.1 to 14.3 by hand.
    \item Bainbridge, ``Ironies of Automation,'' Automatica, 1983.
    \item Sarter and Woods, ``How in the World Did We Ever Get into That Mode?'' Human Factors, 1995.
    \item Zhang et al., ``Determinants of Take-over Time from Automated Driving,'' Transportation Research Part F, 2019.
    \item Sheridan, Telerobotics, Automation, and Human Supervisory Control, MIT Press, 1992.
    \item Ross, Gordon, and Bagnell, ``A Reduction of Imitation Learning and Structured Prediction to No-Regret Online Learning,'' AISTATS 2011.
  \end{itemize}
  \note{[0.5 min] [CORE] The chapter first, with its three checkpoints. Bainbridge explains why the person in reserve is weakest when needed. Sarter and Woods trace mode confusion. Zhang et al. supply the take-over data the deck replotted. Sheridan is the reference on supervisory control. Ross et al. introduce DAgger, the reason the record keeps drift states.}
\end{frame}

\begin{frame}[c]
  \begin{columns}[c]
    \begin{column}{0.34\textwidth}
      \centering\safeimg[width=\textwidth,height=4.2cm,keepaspectratio]{thumb_stack_governance.pdf}
    \end{column}
    \begin{column}{0.62\textwidth}
      {\small\textcolor{lightgray}{NEXT LECTURE \textbar{} ADVERSARIAL VERIFICATION}}\\[0.5cm]
      {\LARGE\bfseries Why can a machine pass its normal tests and still fail when its protection is needed?\par}
    \end{column}
  \end{columns}
  \note{[0.5 min] [CORE] Weeks of clean handovers in the authority record rarely contain a duplicate writer, a replayed packet, or a dropped heartbeat, and a record cannot testify about a fault that never occurred. Lecture 15 injects those faults on purpose, starting with rules A1 to A4 from today's manifest, and scores each trial against an oracle fixed before the run.}
\end{frame}

\end{document}
